% GENERATED FILE. Edit canonical lessons, then run npm run generate:books.
% canonical-source-hash: fnv1a64:425b03395a0f017d
% canonical-lessons: JA-C448-koru, JA-C448-momu, JA-C448-hineru, JA-C448-shibireru, JA-C448-mageru

\chapter{Stiff and Numb}
\label{ch:ja-stiff-and-numb}
\hlchaptermodality{448}

\begin{chapteropening}
\textbf{By the end of this chapter:} \emph{I can say a part of my body is stiff, twisted or numb.}

Complete the last lesson of chapter 448: I can say a part of my body is stiff, twisted or numb.
\end{chapteropening}

\section[\texorpdfstring{\ja{こる} (koru)}{こる (koru)}]{\ja{こる} (koru) — to get stiff (of muscles)}

\label{lesson:JA-C448-koru}



\begin{quote}
Before the new one: say the Japanese for worn out, then the Japanese for lack of sleep.
\end{quote}

\subsection*{You'll want to know: \ja{こる}}
\textbf{\ja{こる}} — \emph{koru} — \textquotedblleft{}to get stiff (of muscles)\textquotedblright{}.

\begin{quote}
\textbf{\ja{しごとで} \ja{かたが} \ja{こりました。}} — \emph{shigoto de kata ga korimashita} — \textquotedblleft{}My shoulders got stiff from work.\textquotedblright{}
\end{quote}

Polite \textbf{\ja{こります}}, \emph{korimasu}; a godan verb. Past (polite): \textbf{\ja{ちちの} \ja{かたが} \ja{こりました。}}, \emph{chichi no kata ga korimashita}.

\begin{cousinweb}
\textbf{\ja{こる}} (\emph{koru}) is a native Japanese word (wago), inherited rather than borrowed from Chinese; its kanji spelling was attached to it later (JMdict).
\end{cousinweb}

\begin{grammarlens}[title={What you've built}]
One more word for describing what you lived and felt.
\end{grammarlens}

\subsection*{Your turn}
\begin{itemize}
\raggedright
  \item \emph{Say it:} \emph{koru}
  \item \emph{Say it:} \emph{koru}, once more
  \item \emph{Say it:} \emph{shigoto de kata ga korimashita} — My shoulders got stiff from work.
  \item \emph{Recall:} say the Japanese for worn out, then the Japanese for lack of sleep, then say \emph{koru} again
\end{itemize}

\begin{tcolorbox}[breakable,colback=teal!4,colframe=teal!35!black,title={Before you move on}]
What does \ja{こる} mean? (\textquotedblleft{}To get stiff (of muscles)\textquotedblright{}.) Say it once more.
\end{tcolorbox}
\section[\texorpdfstring{\ja{もむ} (momu)}{もむ (momu)}]{\ja{もむ} (momu) — to massage; to rub}

\label{lesson:JA-C448-momu}



\begin{quote}
Before the new one: say the Japanese for lack of sleep, then the Japanese for to get stiff.
\end{quote}

\subsection*{You'll want to know: \ja{もむ}}
\textbf{\ja{もむ}} — \emph{momu} — \textquotedblleft{}to massage; to rub\textquotedblright{}.

\begin{quote}
\textbf{\ja{かたを} \ja{もんで} \ja{ください。}} — \emph{kata o monde kudasai} — \textquotedblleft{}Please massage my shoulders.\textquotedblright{}
\end{quote}

Polite \textbf{\ja{もみます}}, \emph{momimasu}; a godan verb. Past (polite): \textbf{\ja{むすめが} \ja{そぼの} \ja{かたを} \ja{もみました。}}, \emph{musume ga sobo no kata o momimashita}.

\begin{cousinweb}
\textbf{\ja{もむ}} (\emph{momu}) is a native Japanese word (wago), inherited rather than borrowed from Chinese; its kanji spelling was attached to it later (JMdict).
\end{cousinweb}

\begin{grammarlens}[title={What you've built}]
One more word for describing what you lived and felt.
\end{grammarlens}

\subsection*{Your turn}
\begin{itemize}
\raggedright
  \item \emph{Say it:} \emph{momu}
  \item \emph{Say it:} \emph{momu}, once more
  \item \emph{Say it:} \emph{kata o monde kudasai} — Please massage my shoulders.
  \item \emph{Recall:} say the Japanese for lack of sleep, then the Japanese for to get stiff, then say \emph{momu} again
\end{itemize}

\begin{tcolorbox}[breakable,colback=teal!4,colframe=teal!35!black,title={Before you move on}]
What does \ja{もむ} mean? (\textquotedblleft{}To massage; to rub\textquotedblright{}.) Say it once more.
\end{tcolorbox}
\section[\texorpdfstring{\ja{ひねる} (hineru)}{ひねる (hineru)}]{\ja{ひねる} (hineru) — to twist}

\label{lesson:JA-C448-hineru}



\begin{quote}
Before the new one: say the Japanese for to get stiff, then the Japanese for to massage.
\end{quote}

\subsection*{You'll want to know: \ja{ひねる}}
\textbf{\ja{ひねる}} — \emph{hineru} — \textquotedblleft{}to twist\textquotedblright{}.

\begin{quote}
\textbf{\ja{ころんで、あしくびを} \ja{ひねりました。}} — \emph{koronde, ashikubi o hinerimashita} — \textquotedblleft{}I fell over and twisted my ankle.\textquotedblright{}
\end{quote}

Polite \textbf{\ja{ひねります}}, \emph{hinerimasu}; a godan verb. Past (polite): \textbf{\ja{せんしゅが} \ja{こしを} \ja{ひねりました。}}, \emph{senshu ga koshi o hinerimashita}.

\begin{cousinweb}
\textbf{\ja{ひねる}} (\emph{hineru}) is a native Japanese word (wago), inherited rather than borrowed from Chinese; its kanji spelling was attached to it later (JMdict).
\end{cousinweb}

\begin{grammarlens}[title={What you've built}]
One more word for describing what you lived and felt.
\end{grammarlens}

\subsection*{Your turn}
\begin{itemize}
\raggedright
  \item \emph{Say it:} \emph{hineru}
  \item \emph{Say it:} \emph{hineru}, once more
  \item \emph{Say it:} \emph{koronde, ashikubi o hinerimashita} — I fell over and twisted my ankle.
  \item \emph{Recall:} say the Japanese for to get stiff, then the Japanese for to massage, then say \emph{hineru} again
\end{itemize}

\begin{tcolorbox}[breakable,colback=teal!4,colframe=teal!35!black,title={Before you move on}]
What does \ja{ひねる} mean? (\textquotedblleft{}To twist\textquotedblright{}.) Say it once more.
\end{tcolorbox}
\section[\texorpdfstring{\ja{しびれる} (shibireru)}{しびれる (shibireru)}]{\ja{しびれる} (shibireru) — to go numb}

\label{lesson:JA-C448-shibireru}



\begin{quote}
Before the new one: say the Japanese for to massage, then the Japanese for to twist.
\end{quote}

\subsection*{You'll want to know: \ja{しびれる}}
\textbf{\ja{しびれる}} — \emph{shibireru} — \textquotedblleft{}to go numb\textquotedblright{}.

\begin{quote}
\textbf{\ja{たたみの} \ja{うえで、あしが} \ja{しびれました。}} — \emph{tatami no ue de, ashi ga shibiremashita} — \textquotedblleft{}My legs went numb on the tatami.\textquotedblright{}
\end{quote}

Polite \textbf{\ja{しびれます}}, \emph{shibiremasu}; an ichidan verb. Past (polite): \textbf{\ja{おとうとの} \ja{あしが} \ja{しびれました。}}, \emph{otōto no ashi ga shibiremashita}.

\begin{cousinweb}
\textbf{\ja{しびれる}} (\emph{shibireru}) is a native Japanese word (wago), inherited rather than borrowed from Chinese; its kanji spelling was attached to it later (JMdict).
\end{cousinweb}

\begin{grammarlens}[title={What you've built}]
One more word for describing what you lived and felt.
\end{grammarlens}

\subsection*{Your turn}
\begin{itemize}
\raggedright
  \item \emph{Say it:} \emph{shibireru}
  \item \emph{Say it:} \emph{shibireru}, once more
  \item \emph{Say it:} \emph{tatami no ue de, ashi ga shibiremashita} — My legs went numb on the tatami.
  \item \emph{Recall:} say the Japanese for to massage, then the Japanese for to twist, then say \emph{shibireru} again
\end{itemize}

\begin{tcolorbox}[breakable,colback=teal!4,colframe=teal!35!black,title={Before you move on}]
What does \ja{しびれる} mean? (\textquotedblleft{}To go numb\textquotedblright{}.) Say it once more.
\end{tcolorbox}
\section[\texorpdfstring{\ja{まげる} (mageru)}{まげる (mageru)}]{\ja{まげる} (mageru) — to bend}

\label{lesson:JA-C448-mageru}



\begin{quote}
Before the new one: say the Japanese for to twist, then the Japanese for to go numb.
\end{quote}

\subsection*{You'll want to know: \ja{まげる}}
\textbf{\ja{まげる}} — \emph{mageru} — \textquotedblleft{}to bend\textquotedblright{}.

\begin{quote}
\textbf{\ja{ひざを} \ja{ゆっくり} \ja{まげて} \ja{ください。}} — \emph{hiza o yukkuri magete kudasai} — \textquotedblleft{}Please bend your knees slowly.\textquotedblright{}
\end{quote}

Polite \textbf{\ja{まげます}}, \emph{magemasu}; an ichidan verb. Past (polite): \textbf{\ja{かんじゃが} \ja{ゆびを} \ja{まげました。}}, \emph{kanja ga yubi o magemashita}.

\begin{cousinweb}
\textbf{\ja{まげる}} (\emph{mageru}) is a native Japanese word (wago), inherited rather than borrowed from Chinese: the transitive partner of \emph{magaru}, \textquotedblleft{}to turn; to bend\textquotedblright{}, and its kanji was attached later as a kun'yomi (JMdict, KANJIDIC2).
\end{cousinweb}

\begin{grammarlens}[title={What you've built}]
One more word for describing what you lived and felt.
\end{grammarlens}

\subsection*{Your turn}
\begin{itemize}
\raggedright
  \item \emph{Say it:} \emph{mageru}
  \item \emph{Say it:} \emph{mageru}, once more
  \item \emph{Say it:} \emph{hiza o yukkuri magete kudasai} — Please bend your knees slowly.
  \item \emph{Recall:} say the Japanese for to twist, then the Japanese for to go numb, then say \emph{mageru} again
\end{itemize}

\begin{tcolorbox}[breakable,colback=teal!4,colframe=teal!35!black,title={Before you move on}]
What does \ja{まげる} mean? (\textquotedblleft{}To bend\textquotedblright{}.) Say it once more.
\end{tcolorbox}
